% registry: SIH
\methodhead{SIH}{Zenner, Heidenreich and Richter: shear stress intensity hypothesis}\label{sheet:SIH}

\block{Source} H.~Zenner, A.~Simb\"urger, J.~Liu, \emph{Int.\ J.\ Fatigue} 22 (2000) 137--145
\cite{zenner2000}, Eqs.~(4)--(7), Table~1. The hypothesis is that of Zenner, Heidenreich and
Richter (1983); the shear amplitude for a rotating shear vector is defined in Liu (1991) and Liu
and Zenner (1993), which were not consulted.

\block{Idea and scope} Integral approach. The squared shear and normal stress amplitudes, each
corrected by its mean stress, are averaged over all planes through the point. The mean square of
the shear stress over all planes is proportional to $J_2'$ (Novozhilov, Eq.~4). So for constant
principal directions the hypothesis reduces to the octahedral criterion. When the principal
directions rotate, the shear amplitude is spread over more planes instead of being lowered on the
worst one. Proposed for fatigue limits of ductile steels ($R_m$ = 400--1600\MPa) with
$0.577\le\tau_W/\sigma_W\le0.866$. The source covers alternating and static normal and shear
stresses, phase shifts, different frequencies and non-sinusoidal waveforms (179 test series,
Table~2).

\block{Inputs}
\begin{center}\small
\begin{tabular}{@{}llp{100mm}@{}}
\toprule
Symbol & Unit & Meaning\\\midrule
$\gamma,\varphi$ & rad & polar and azimuth angles of the plane normal, $\gamma\in[0,\pi]$, $\varphi\in[0,2\pi)$\\
$\tau_{\gamma\varphi,a}$, $\tau_{\gamma\varphi,m}$ & MPa & shear stress amplitude and mean on the plane\\
$\sigma_{\gamma\varphi,a}$, $\sigma_{\gamma\varphi,m}$ & MPa & normal stress amplitude and mean on the plane\\
$\sigma_W=\fm$, $\tau_W=\tm$ & MPa & fully reversed axial and torsion limits\\
$\sigma_{Sch}$, $\tau_{Sch}$ & MPa & pulsating ($R=0$) axial and torsion limits, maximum stress\\
$a,b$ & -- & weights of the shear and normal terms\\
$m$, $n$ & MPa$^{-2}$, MPa$^{-1}$ & mean shear and mean normal stress sensitivities\\
$\sigma_{eq,a}$ & MPa & equivalent stress amplitude, Eq.~\eqref{eq:sih}\\
\bottomrule
\end{tabular}
\end{center}

\block{Formulation} The equivalent stress amplitude (Eq.~6) and the failure condition are
\begin{equation}
\sigma_{eq,a}=\left\{\frac{15}{8\pi}\int_0^{\pi}\!\!\int_0^{2\pi}\Bigl[a\,\tau_{\gamma\varphi,a}^2\bigl(1+m\,\tau_{\gamma\varphi,m}^2\bigr)
+b\,\sigma_{\gamma\varphi,a}^2\bigl(1+n\,\sigma_{\gamma\varphi,m}\bigr)\Bigr]\sin\gamma\,d\gamma\,d\varphi\right\}^{1/2}=\sigma_W .
\label{eq:sih}
\end{equation}
The constants (Table~1) make the four limits exact:
\begin{gather}
a=\tfrac15\Bigl[3\Bigl(\frac{\sigma_W}{\tau_W}\Bigr)^2-4\Bigr],\qquad
b=\tfrac15\Bigl[6-2\Bigl(\frac{\sigma_W}{\tau_W}\Bigr)^2\Bigr],\\
a\,m=\frac{\sigma_W^2-(\sigma_W/\tau_W)^2(\tau_{Sch}/2)^2}{\tfrac{12}{7}(\tau_{Sch}/2)^4},\qquad
b\,n=\frac{\sigma_W^2-(\sigma_{Sch}/2)^2-\tfrac{4}{21}\,a\,m\,(\sigma_{Sch}/2)^4}{\tfrac{15}{14}(\sigma_{Sch}/2)^3}.
\end{gather}
Here $12/7$, $4/21$ and $15/14$ are $\tfrac{15}{8\pi}\oint$ of $\tau^4$ (torsion), $\tau^4$
(tension) and $\sigma^3$ (tension) for unit stresses. When $\tau_{Sch}$ is not measured (Eq.~7),
\begin{equation}
\tau_{Sch}=\frac{4\,\tau_W}{2\sigma_W/\sigma_{Sch}+1}.
\end{equation}
With $a=1$, $b=0$ and no means, Eq.~\eqref{eq:sih} is Eq.~(5), which gives $\tau_W/\sigma_W=1/\sqrt3$;
for constant principal directions it equals $\tfrac{3}{\sqrt2}\tau_{oct,a}$ (Eq.~5b).

\block{Constants and calibration} $a,b,m,n$ from $\sigma_W,\tau_W,\sigma_{Sch},\tau_{Sch}$; nothing
is fitted. The source says $\sigma_{Sch}$ may be estimated from the mean stress sensitivity, but
gives no formula.

\block{Life} None: a fatigue-limit criterion. Section~6 states that lives under variable amplitude
need damage sums integrated per plane. The repository reads finite lives on the master curve of
the index $\sigma_{eq,a}/\sigma_W$ (Section~\ref{sec:mastercurve}).

\block{Repository choices}
\begin{itemize}
\item Shear amplitude of a rotating shear vector: half the longest chord of the single-harmonic
  ellipse, i.e.\ its minimum circumscribed circle radius (maintainer decision; the source
  definition is not available). The quadratic amplitude $\sqrt{s^2+c^2}$ is selectable. Section~3
  states that Eq.~(5) does not depend on the phase between $\sigma_x$ and $\tau_{xy}$. That holds
  exactly with the quadratic amplitude; with the longest chord, Eq.~(5) drops by 7\% at
  $90^\circ$ for $\tau_{xy,a}=\sigma_{x,a}/2$. The maintainer confirmed the longest chord with
  this evidence.
\item $\tau_{\gamma\varphi,m}$ is the magnitude of the mean shear vector.
\item Mean stresses require \texttt{Sig0}; \texttt{Tau0} defaults to Eq.~(7).
\item $\tau_W/\sigma_W$ outside $[0.577,0.866]$ is out of material scope. A negative integral is
  clipped to zero.
\item 32$\times$64 Gauss--Legendre $\times$ azimuth quadrature; the model returns
  $\sigma_{eq,a}/\sigma_W$.
\end{itemize}

\block{Worked example} Hard steel of Nishihara and Kawamoto, $\sigma_W=313.9$,
$\tau_W=196.2\MPa$ ($\tau_W/\sigma_W=0.625$): $a=0.7358$, $b=0.1761$. Bending--torsion limits of
Papadopoulos et al.\ (1997), Table~1. In phase, $\sigma_a=138.1$, $\tau_a=167.1\MPa$: index 0.959
with either amplitude. At $90^\circ$, $\sigma_a=150.2$, $\tau_a=181.7\MPa$: index 0.998 with the
longest chord, 1.042 with the quadratic amplitude.

\block{Verification} Checked: the Table~1 factors as sphere integrals; the four limits exact; Eq.~(5)
and Eq.~(5b); the Section~3 phase independence (quadratic amplitude). The paper gives only the
statistics of 179 series (Table~2), not the series.
